\chapter{Introducción}

El envejecimiento de la población es uno de los fenómenos más significativos de la Argentina contemporánea. El aumento de la esperanza de vida y las mejoras sostenidas en las condiciones de salud han permitido que hoy más de 7,4 millones de personas tengan 60 años o más, representando el 15,7\% de la población total, y las proyecciones indican que esa proporción superará el 20\% hacia 2040 \parencite{INDEC2024}. Vivir más años y con mayor autonomía es el resultado de décadas de avances en salud pública, nutrición y calidad de vida.

Sin embargo, este cambio demográfico trae consigo nuevas realidades familiares y sociales que aún no encuentran respuesta adecuada. Solo el 46\% de las personas de 60 años y más vive en hogares multigeneracionales \parencite{INDECTablero2024}, lo que implica que la mayoría enfrenta situaciones de supervisión parcial o nula, y quienes conviven con sus hijos o parientes tampoco cuentan necesariamente con supervisión continua, los hijos que cuidan a sus padres mayores son, en la mayoría de los casos, personas de mediana edad con sus propias familias y trabajos. Esta situación genera una brecha de información cotidiana muy difícil de cerrar. Los familiares no saben si el adulto mayor comió, si se levantó a la hora habitual, o si algo cambió en su rutina.

Esta incertidumbre no es menor si se considera que los adultos mayores que viven solos o sin supervisión durante el día están expuestos a riesgos concretos y silenciosos. El deterioro de la salud muchas veces comienza a través de cambios sutiles en la rutina diaria ---levantarse más tarde, reducir la actividad, saltarse comidas--- que, de ser detectados a tiempo, permitirían actuar antes de que se conviertan en una emergencia. En los casos más críticos, las caídas representan uno de los eventos de mayor impacto: afectan a aproximadamente el 30\% de los adultos mayores independientes cada año \parencite{YeohLui2025}. Sin presencia física o algún mecanismo de observación continua, tanto el deterioro silencioso como los eventos agudos pasan inadvertidos hasta que ya es tarde para actuar.

Las alternativas tecnológicas disponibles hoy en día no resuelven este problema de fondo. Los dispositivos wearables y los botones de pánico requieren que el adulto mayor los use activamente y de forma constante, lo que choca con una realidad documentada: en Argentina, los adultos mayores presentan los niveles más bajos de uso de tecnologías de la información y la comunicación respecto al resto de la población \parencite{INDEC2020}, y esta brecha no responde a limitaciones cognitivas sino a dificultades de uso y preocupaciones sobre la privacidad que dificultan la adopción sostenida de dispositivos \parencite{Lam2025}. El resultado es que las soluciones existentes funcionan bien en teoría, pero tienen bajas tasas de adopción y sostenimiento en la práctica. Frente a este panorama, la alternativa viable es un sistema que opere de forma completamente pasiva, sin requerir ninguna acción por parte del adulto mayor, y es precisamente la visión artificial la tecnología que hace posible ese enfoque.

El presente proyecto propone una respuesta concreta a este problema. Se desarrolla un sistema de monitoreo no invasivo para adultos mayores en entornos domésticos, pensado especialmente para aquellas familias que, por distancia o por las exigencias de su vida cotidiana, no pueden estar presentes junto al adulto mayor a su cargo. El sistema opera de forma completamente pasiva y autónoma: no requiere ninguna intervención por parte del adulto mayor, no utiliza dispositivos que deban portarse, y todo el procesamiento de video se realiza de manera local en el hogar, garantizando que los registros visuales no salgan del domicilio.

El sistema integra dos módulos principales. El primero analiza el video de cámaras IP en tiempo real para detectar eventos críticos como caídas o inmovilidad prolongada. El segundo construye y actualiza de forma continua una línea base de comportamiento personalizada para cada individuo, permitiendo identificar desvíos significativos en la rutina diaria que podrían indicar un deterioro en el estado de salud. Los eventos detectados y las descripciones de actividad generadas se transmiten a la nube para notificar a la familia y alimentar los reportes periódicos a través de una aplicación móvil.

% Párrafo de estructura del informe - completar con las secciones finales acordadas
El presente informe se organiza de la siguiente manera. En la sección de Antecedentes se desarrolla el marco teórico que sustenta el proyecto, el análisis del estado del arte de las soluciones existentes y el user research mediante entrevistas con familiares de adultos mayores, con el objetivo de validar la problemática identificada y orientar las decisiones de diseño del sistema. En la sección de Descripción se especifica el sistema desarrollado, incluyendo su arquitectura, requerimientos y componentes principales. En Metodología de Desarrollo se detalla el proceso de construcción, las tecnologías utilizadas y las decisiones de diseño tomadas. En Pruebas Realizadas se presentan los ensayos efectuados sobre el prototipo funcional. Finalmente, en la sección de Discusión se analizan los resultados obtenidos, las limitaciones encontradas y las diferencias entre el alcance planificado y el real, seguido de las Conclusiones del trabajo.

\section{Objetivos}

\subsection*{Objetivo general}

Desarrollar un sistema de monitoreo no invasivo para adultos mayores en entornos domésticos, basado en visión artificial y aprendizaje adaptativo del comportamiento, que permita a sus familiares mantenerse informados sobre su estado y rutina diaria, y recibir alertas ante situaciones de riesgo, sin requerir ninguna intervención activa por parte del adulto mayor.

\subsection*{Objetivos específicos}

\begin{itemize}
    \item Implementar un módulo de análisis de video en tiempo real capaz de detectar caídas e inmovilidad prolongada a partir de modelos de estimación de poses.
    \item Desarrollar un módulo de clasificación semántica de actividades de la vida diaria mediante consultas periódicas a un modelo de lenguaje visual (VLM), que complemente la inferencia en tiempo real con descripciones textuales de la actividad observada.
    \item Desarrollar un módulo de aprendizaje adaptativo que construya y actualice de forma continua una línea base de comportamiento personalizada por individuo, permitiendo identificar desvíos significativos en la rutina diaria.
    \item Garantizar el procesamiento local del video en el hogar, de modo que los registros visuales no salgan del domicilio, preservando la privacidad del adulto mayor.
    \item Diseñar e implementar una aplicación móvil que notifique a los familiares ante eventos detectados y permita el acceso a reportes periódicos sobre el estado y la rutina del adulto mayor.
    \item Validar el sistema en un entorno doméstico controlado mediante escenarios simulados que permitan evaluar la efectividad de la detección.
\end{itemize}
\vspace*{4pt}

\section{Alcance}

\subsection{Características del prototipo funcional}

El proyecto abarca el diseño e implementación de un prototipo funcional compuesto por dos módulos principales. El primero realiza el análisis visual mediante un pipeline que combina un modelo de estimación de poses para la detección en tiempo real de eventos agudos ---caídas e inmovilidad prolongada--- con consultas periódicas a un modelo de lenguaje visual (VLM) para la clasificación semántica de actividades de la vida diaria (AVD). El segundo construye y actualiza de forma continua una línea base de comportamiento personalizada por usuario mediante clustering dinámico, permitiendo identificar desvíos significativos en la rutina diaria como señales tempranas de deterioro funcional. El procesamiento de video se ejecuta localmente en una controladora conectada a cámaras IP, garantizando que los registros visuales no salgan del domicilio; únicamente los eventos detectados y las descripciones de actividad generadas se transmiten a la nube para el envío de notificaciones y el registro de actividad. El sistema incluye además una aplicación móvil para la visualización de alertas y reportes, y se valida en un entorno doméstico controlado mediante escenarios simulados.

\subsection{Limitaciones y exclusiones}

El prototipo no contempla integración con sistemas médicos externos ni producción a escala comercial. Tampoco incluye soporte para múltiples adultos mayores en distintos hogares de forma simultánea, ni vinculación con servicios de emergencia o atención médica profesional.

\subsection{Estructura del documento}

El presente informe se organiza en seis secciones principales. La sección de Antecedentes desarrolla el marco teórico que sustenta el proyecto, el análisis del estado del arte de las soluciones existentes y el user research mediante entrevistas con familiares de adultos mayores, con el objetivo de validar la problemática identificada y orientar las decisiones de diseño del sistema. La sección de Descripción especifica el sistema desarrollado, incluyendo su arquitectura, requerimientos y componentes principales. En Metodología de Desarrollo se detalla el proceso de construcción, las tecnologías utilizadas y las decisiones de diseño tomadas. En Pruebas Realizadas se presentan los ensayos efectuados sobre el prototipo funcional. Finalmente, en la sección de Discusión se analizan los resultados obtenidos, las limitaciones encontradas y las diferencias entre el alcance planificado y el real, seguido de las Conclusiones del trabajo.